\onecolumn
\section{Detailed Step-by-Step Derivation of the Trajectory Bound}
\label{app:detailed_proof}

\subsection*{Setup and Definitions}

\paragraph{1. The Combined Loss.}
The total loss function at any point $\lambda \in [0,1]$ in the curriculum is defined as:
\begin{equation}
    L_\lambda(\theta) = L_{\text{replay}}(\theta) + \lambda \cdot L_{\text{current}}(\theta)
\end{equation}

\paragraph{2. Defining the Hessians ($\nabla^2$).}
Differentiating the gradient once more gives the second-derivative matrix, the Hessian ($H$).
\begin{itemize}
    \item $H_{\text{replay}} = \nabla^2 L_{\text{replay}}(\theta)$ (The curvature of the old task)
    \item $H_{\text{current}} = \nabla^2 L_{\text{current}}(\theta)$ (The curvature of the new task)
\end{itemize}
Therefore, the Hessian of our combined loss, $H_\lambda$, is defined by applying the second derivative operator to the total loss function:
\begin{align*}
    H_\lambda & = \nabla^2 \left[ L_{\text{replay}}(\theta) + \lambda \cdot L_{\text{current}}(\theta) \right] \\
    H_\lambda & = \nabla^2 L_{\text{replay}}(\theta) + \lambda \cdot \nabla^2 L_{\text{current}}(\theta)       \\
    H_\lambda & = H_{\text{replay}} + \lambda \cdot H_{\text{current}}
\end{align*}

\paragraph{3. Assumption: Local-Minimum along the Interpolation Branch.}
We assume that the loss functions $L_{\text{replay}}$ and $L_{\text{current}}$ are twice continuously differentiable ($C^2$) in a neighbourhood of the interpolation path. Furthermore, we assume there exists a continuously differentiable ($C^1$) branch of critical points $\lambda \mapsto \Theta^*(\lambda)$ for the interpolated loss $L_\lambda$, initialized at $\Theta^*(0) = \theta_0^*$, such that the Hessian remains strictly positive-definite: $H_\lambda(\Theta^*(\lambda)) \succ 0$ for every $\lambda \in [0,1]$.

This assumption ensures that the iterate tracks a continuous family of strict local minima rather than relying on global convexity. Near $\lambda = 0$, this condition is naturally satisfied. Because the initial state $\theta_0^*$ is a non-degenerate local minimum of $L_{\text{replay}}$, it follows that $H_0 = H_{\text{replay}}(\theta_0^*) \succ 0$. Consequently, the implicit function theorem guarantees a unique smooth branch in the immediate neighbourhood of $\lambda = 0$.

The persistence of this positive-definite branch as $\lambda$ increases is bounded by Weyl's inequality applied to $H_\lambda = H_{\text{replay}} + \lambda H_{\text{current}}$; writing $\mu_{\min}(\cdot)$ for the smallest eigenvalue,
\begin{equation*}
    \mu_{\min}(H_\lambda) \;\ge\; \mu_{\min}(H_{\text{replay}}) + \lambda\, \mu_{\min}(H_{\text{current}}).
\end{equation*}
If the new task is locally convex at these points ($H_{\text{current}} \succeq 0$), positive-definiteness holds for all $\lambda \ge 0$ without requiring a global convexity assumption.

\subsection*{Part A: The Replay Loss Is Monotone Along the Reference Path}

Let $\Theta^*(\lambda)$ be the reference path, the exact parameter state that minimises $L_\lambda(\theta)$ for any given $\lambda$.

\paragraph{Step 1: The First-Order Condition.}
Because $\Theta^*(\lambda)$ is a minimum, the gradient of the total loss evaluated at that exact point must be zero:
\begin{align*}
    \nabla \left[ L_\lambda(\Theta^*(\lambda)) \right]                                                               & = 0 \\
    \nabla \left[ L_{\text{replay}}(\Theta^*(\lambda)) + \lambda \cdot L_{\text{current}}(\Theta^*(\lambda)) \right] & = 0
\end{align*}
Distributing the gradient operator to both terms gives us the fundamental First-Order Condition:
\begin{equation}
    \nabla L_{\text{replay}}(\Theta^*(\lambda)) + \lambda \cdot \nabla L_{\text{current}}(\Theta^*(\lambda)) = 0
    \label{eq:foc_detailed}
\end{equation}

\paragraph{Step 2: Implicit Differentiation to Find Path Movement.}
We want to know how the reference path $\Theta^*$ moves as $\lambda$ changes. We apply the derivative operator $\frac{d}{d\lambda}$ to the entire First-Order Condition:
\begin{equation*}
    \frac{d}{d\lambda} \left[ \nabla L_{\text{replay}}(\Theta^*(\lambda)) + \lambda \cdot \nabla L_{\text{current}}(\Theta^*(\lambda)) \right] = \frac{d}{d\lambda} [0]
\end{equation*}
Distributing the operator to each term:
\begin{equation*}
    \frac{d}{d\lambda} \left[ \nabla L_{\text{replay}}(\Theta^*(\lambda)) \right] + \frac{d}{d\lambda} \left[ \lambda \cdot \nabla L_{\text{current}}(\Theta^*(\lambda)) \right] = 0
\end{equation*}
Evaluating the first term via the multivariable chain rule:
\begin{equation*}
    \frac{d}{d\lambda} \left[ \nabla L_{\text{replay}}(\Theta^*(\lambda)) \right] = \nabla^2 L_{\text{replay}}(\Theta^*(\lambda)) \cdot \frac{d\Theta^*}{d\lambda}
\end{equation*}
Evaluating the second term via the product rule and chain rule:
\begin{equation*}
    \frac{d}{d\lambda} \left[ \lambda \cdot \nabla L_{\text{current}}(\Theta^*(\lambda)) \right] = (1) \cdot \nabla L_{\text{current}}(\Theta^*(\lambda)) + \lambda \cdot \left( \nabla^2 L_{\text{current}}(\Theta^*(\lambda)) \cdot \frac{d\Theta^*}{d\lambda} \right)
\end{equation*}
Substituting these expanded terms back into the equation:
\begin{equation*}
    \left[ \nabla^2 L_{\text{replay}}(\Theta^*) \cdot \frac{d\Theta^*}{d\lambda} \right] + \left[ \nabla L_{\text{current}}(\Theta^*) + \lambda \cdot \nabla^2 L_{\text{current}}(\Theta^*) \cdot \frac{d\Theta^*}{d\lambda} \right] = 0
\end{equation*}
Grouping the terms that share the path movement $\frac{d\Theta^*}{d\lambda}$:
\begin{equation*}
    \left( \nabla^2 L_{\text{replay}}(\Theta^*) + \lambda \cdot \nabla^2 L_{\text{current}}(\Theta^*) \right) \cdot \frac{d\Theta^*}{d\lambda} + \nabla L_{\text{current}}(\Theta^*) = 0
\end{equation*}
Recognising the definition of the combined Hessian $H_\lambda$ inside the parentheses:
\begin{equation*}
    H_\lambda \cdot \frac{d\Theta^*}{d\lambda} + \nabla L_{\text{current}}(\Theta^*(\lambda)) = 0
\end{equation*}
Solving algebraically for the path movement $\frac{d\Theta^*}{d\lambda}$:
\begin{equation}
    \frac{d\Theta^*}{d\lambda} = -H_\lambda^{-1} \cdot \nabla L_{\text{current}}(\Theta^*(\lambda))
    \label{eq:path_movement}
\end{equation}

\paragraph{Step 3: Calculating the Change in the Replay Loss.}
How does the replay loss change as $\lambda$ increases along this path? We apply the derivative operator:
\begin{equation*}
    \frac{d}{d\lambda} \left[ L_{\text{replay}}(\Theta^*(\lambda)) \right] = \nabla L_{\text{replay}}(\Theta^*(\lambda))^\top \cdot \frac{d\Theta^*(\lambda)}{d\lambda}
\end{equation*}
We perform two substitutions here. First, from \eqref{eq:foc_detailed}, we know $\nabla L_{\text{replay}} = -\lambda \cdot \nabla L_{\text{current}}$. Second, from \eqref{eq:path_movement}, we substitute the path movement.
\begin{align*}
    \frac{d}{d\lambda} L_{\text{replay}}(\Theta^*(\lambda)) & = \left( -\lambda \cdot \nabla L_{\text{current}}(\Theta^*(\lambda)) \right)^\top \cdot \left( -H_\lambda^{-1} \cdot \nabla L_{\text{current}}(\Theta^*(\lambda)) \right) \\
                                                            & = \lambda \cdot \nabla L_{\text{current}}(\Theta^*(\lambda))^\top \cdot H_\lambda^{-1} \cdot \nabla L_{\text{current}}(\Theta^*(\lambda))
\end{align*}
Because $H_\lambda$ is strictly positive-definite along the followed branch of local minima for all $\lambda \in [0,1]$ (per the Local-Minimum Assumption), its inverse $H_\lambda^{-1}$ is also positive-definite. The quadratic form $v^\top M v$ is therefore guaranteed to be non-negative. Since $\lambda \ge 0$, the entire equation satisfies:
\begin{equation*}
    \frac{d}{d\lambda} L_{\text{replay}}(\Theta^*(\lambda)) \ge 0
\end{equation*}
\emph{Conclusion for Part A:} The replay loss never decreases along the reference path, so it can never overshoot its final joint-optimum value.

\subsection*{Part B: The Iterate Tracks the Reference Path with $O(1/N)$ Lag}

Let $\theta(t)$ be the actual position of the model parameters at time $t$. We define the tracking error $e(t)$ (written $e$ to avoid a clash with the damping $\delta$ of asymmetric PER):
\begin{equation*}
    e(t) = \theta(t) - \Theta^*(\lambda(t))
\end{equation*}

\paragraph{Step 1: Differentiating the Error.}
We apply the time derivative operator $\frac{d}{dt}$ to our error definition:
\begin{equation*}
    \frac{d}{dt} \left[ e(t) \right] = \frac{d}{dt} \left[ \theta(t) \right] - \frac{d}{dt} \left[ \Theta^*(\lambda(t)) \right]
\end{equation*}
The model moves via gradient descent, so $\frac{d\theta}{dt} = -\nabla L_{\lambda(t)}(\theta(t))$. The reference path moves via the chain rule $\frac{d\Theta^*}{dt} = \frac{d\Theta^*}{d\lambda} \cdot \frac{d\lambda}{dt}$. Because the curriculum has a constant velocity over $N$ steps, $\frac{d\lambda}{dt} = \frac{1}{N}$.
\begin{equation*}
    \frac{de}{dt} = -\nabla L_{\lambda(t)}(\theta(t)) - \left( \frac{d\Theta^*}{d\lambda} \cdot \frac{1}{N} \right)
\end{equation*}

\paragraph{Step 2: Taylor Expansion of the Actual Gradient.}
We approximate the gradient at the actual position $\theta(t)$ by Taylor-expanding it around the reference position $\Theta^*$:
\begin{equation*}
    \nabla L_\lambda(\theta) \approx \nabla L_\lambda(\Theta^*) + \nabla \left[ \nabla L_\lambda(\Theta^*) \right] \cdot (\theta - \Theta^*)
\end{equation*}
Since the gradient at the minimum is zero ($\nabla L_\lambda(\Theta^*) = 0$) and the derivative of the gradient is the Hessian ($\nabla^2 L_\lambda = H_\lambda$), this collapses to:
\begin{equation*}
    \nabla L_\lambda(\theta) \approx H_\lambda \cdot e
\end{equation*}

\paragraph{Step 3: The Final Steady State.}
Substitute this expanded gradient back into our differential equation:
\begin{equation*}
    \frac{de}{dt} \approx -H_\lambda \cdot e - \frac{d\Theta^*}{d\lambda} \cdot \frac{1}{N}
\end{equation*}
Assuming the lag reaches a steady state, the error stops growing ($\frac{de}{dt} \approx 0$). This is a first-order, gradient-flow approximation: it linearises the gradient around $\Theta^*$ and treats the dynamics as continuous, so it captures the scaling of the lag with $N$ rather than an exact discrete-time bound.
\begin{align*}
    0                & \approx -H_\lambda \cdot e - \frac{d\Theta^*}{d\lambda} \cdot \frac{1}{N} \\
    H_\lambda \cdot e & \approx -\frac{d\Theta^*}{d\lambda} \cdot \frac{1}{N}                          \\
    e                & \approx -H_\lambda^{-1} \cdot \frac{d\Theta^*}{d\lambda} \cdot \frac{1}{N}
\end{align*}
Taking the norm, we see the lag distance is bounded by the curriculum length $N$:
\begin{equation}
    \|e\| = O\left(\frac{1}{N}\right)
\end{equation}

\subsection*{Combining Part A and B: Expanding the Actual Loss}

To find the maximum replay loss the model will actually experience, we Taylor expand the loss of the actual trajectory around the reference position:
\begin{align*}
    L_{\text{replay}}(\theta) & \approx L_{\text{replay}}(\Theta^*) + \nabla L_{\text{replay}}(\Theta^*)^\top \cdot (\theta - \Theta^*) \\
    L_{\text{replay}}(\theta) & \approx L_{\text{replay}}(\Theta^*) + \nabla L_{\text{replay}}(\Theta^*)^\top \cdot e
\end{align*}
From Part A, $L_{\text{replay}}(\Theta^*)$ cannot exceed the final joint optimum loss $L_{\text{replay}}(\theta_{\text{joint}}^*)$. From Part B, the lag $e$ is $O(1/N)$. Substituting these limits yields the final bound:
\begin{equation}
    \max_t L_{\text{replay}}(\theta(t)) \le L_{\text{replay}}(\theta_{\text{joint}}^*) + O\left(\frac{1}{N}\right)
\end{equation}
\emph{Conclusion:} The transient stability gap is entirely contained within the $O(1/N)$ term, so a sufficiently slow curriculum makes the tracking lag's contribution to the transient arbitrarily small.

